\section[Wrap-up]{Wrap-up}

\begin{frame}
    \frametitle{You wrote all of this}
    \vfill
    \begin{roadmap}{5}
        \roadmapresults
    \end{roadmap}
    \note{%
        Two hours ago, most of them had never compiled a document.
        The whole line is done.
        Click through the four lists slowly, one block at a time.
        \begin{itemize}
            \item This is the slide the advance organiser was built for:
                four small things are now visibly one finished thing.
                Do not talk over the last click.
        \end{itemize}
    }
\end{frame}

\begin{frame}
    \frametitle{Three packages to add next}
    \begin{itemize}
        \item<1-> \textbf{\texttt{siunitx}}: numbers with units
            \tcbinputlst[linenos=false][left=2mm, overlay={}]{listings/99-siunitx.tex}
            \begin{outputbox}[width=.6\textwidth]
                \qty{65}{\degreeCelsius}, \qty{300}{\milli\litre}, \qty{24}{\minute}

            \end{outputbox}
        \item<2-> \textbf{\texttt{microtype}}: better justification
        \item<3-> \textbf{\texttt{hyperref}}: clickable references, load it \textbf{last}
    \end{itemize}
    \note{%
        \begin{itemize}
            \item Three lines in the preamble, no change to anything they have written.
            \item \texttt{siunitx}: right spacing, right font, every time.
            \item \texttt{microtype}: load it and look closely at the right margin.
            \item \texttt{hyperref} makes every \texttt{\textbackslash ref}, \texttt{\textbackslash eqref} and \texttt{\textbackslash cite} clickable.
            \item \texttt{siunitx} comes first because you owed it to them since block 3,
                where they typed \texttt{\textbackslash,\^{}\{\textbackslash circ\}\textbackslash mathrm\{C\}} by hand.
                Close the loop: ``remember that mess? This.''
            \item Deliberately only three.
                A list of twenty useful packages is a list nobody acts on.
        \end{itemize}
    }
\end{frame}

\begin{frame}
    \frametitle{When your document gets bigger}
    \begin{columns}[T, onlytextwidth]
        \begin{column}{.4\textwidth}
            \tcbinputlst[linenos=false][minted language=text, left=2mm, overlay={}]{listings/99-tree.txt}
        \end{column}
        \begin{column}{.52\textwidth}
            \tcbinputlst{listings/99-input.tex}
            One file per chapter
        \end{column}
    \end{columns}
    \note{%
        Thirty seconds.
        Not worth it for three pages, essential at thirty:
        the editor stays usable, and two people can work without fighting.
        They do not need this today and will need it for their thesis.
        \begin{itemize}
            \item The point is that they have seen it once
                and recognise it when a supervisor's template does it.
            \item Same for a separate \texttt{preamble.tex}: mention it, do not demo it.
        \end{itemize}
    }
\end{frame}

\begin{frame}
    \frametitle{Where to look things up}
    \begin{columns}[T, onlytextwidth]
        \begin{column}{.52\textwidth}
            \begin{itemize}
                \item \textbf{Detexify}: symbol to command\\
                    {\small\color{black!60}detexify.kirelabs.org}
                \item \textbf{Overleaf docs}: beginner reference\\
                    {\small\color{black!60}overleaf.com/learn}
                \item \textbf{TeX StackExchange}: errors, answered\\
                    {\small\color{black!60}tex.stackexchange.com}
                \item \textbf{CTAN}: packages and their manuals\\
                    {\small\color{black!60}ctan.org}
            \end{itemize}
        \end{column}
        \begin{column}{.44\textwidth}
            \begin{block}{Also in the handout}
                \begin{itemize}
                    \item all the sectioning levels
                    \item the rest of the text styling commands
                    \item matrix environments
                    \item column types and float specifiers
                    \item the reserved characters, in full
                \end{itemize}
            \end{block}
        \end{column}
    \end{columns}
    \note{%
        The handout is written in \LaTeX{}; its source is in the repository.
        Hand the handout out at the \emph{start} of the session, not here.
        It is a reference they should have in front of them all afternoon.
        This slide is just the reminder that it exists, and why the deck left things out.
    }
\end{frame}

\begin{frame}
    \frametitle{One last thing}
    \vfill
    \begin{center}
        {\LARGE Nobody writes \LaTeX{} from memory\par}
        \vspace{1cm}
        \uncover<2->{%
            {\Large\color{black!60}
            Look it up, copy your last preamble\par}
        }
    \end{center}
    \vfill
    \note{%
        The honest closing note, worth ten seconds.
        The beginner's failure mode is believing that fluent users have memorised it all
        and that they are uniquely slow.
        They have not, and the student is not.
    }
\end{frame}

\begin{frame}
    \frametitle{Thank you}
    \begin{columns}[c, onlytextwidth]
        \begin{column}{.55\textwidth}
            {\Large Template, slides and handout\par}
            \medskip
            {\small\url{\repourl}\par}
            \vspace{1cm}
            {\Large Questions?\par}
        \end{column}
        \begin{column}{.35\textwidth}
            \centering
            \qrcode[level=L, version=0, height=4cm]{\repourl}
        \end{column}
    \end{columns}
    \note{%
        If something breaks in their next report, they can ask.
        Stay at the front for a few minutes if you can.
        The best questions come one-to-one after the room empties.
    }
\end{frame}

% TODO before publishing: fill in the original author and link (also in the
% repository's top-level LICENSE), or move this text to the README and delete
% the frame.
\begin{frame}
    \frametitle{Attribution and licence}
    \begin{center}
        \begin{minipage}{.8\textwidth}
            \textbf{Based on}:
            \textcolor{psecbg}{TODO: original author}, \textcolor{psecbg}{TODO: link}, CC~BY-SA~4.0

            \bigskip
            \textbf{This revision}: Laurin Heitzer, ECE, \the\year, CC~BY-SA~4.0

            \bigskip
            {\small\color{black!60}
            Example document and figure: written for this workshop, synthetic data\par}
        \end{minipage}
    \end{center}
    \note{%
        For you, not the students.
        Either fill it in and keep it as a visible credit slide,
        or move the text to the README and delete the slide.
        Share-alike is not optional: a CC~BY-SA original makes this deck CC~BY-SA too.
    }
\end{frame}
